\section{怎样辨析单句和复句}
复句是由两个或两个以上的单句组成的一个较大的预言单位。组成复句的单句叫分句，分句可以是主谓句，也可以是非主谓句。例如：

（1）我们不但善于破坏一个旧世界，我们还将善于建设一个新世界。

（2）我们能够去掉不良作风，保持优良作风。

（3）没有四个现代化，社会主义对资本主义的最后胜利是不可能的。

（4）下了一阵急雨，竟然出起太阳来了。

第一句的两个分句是主谓句，都有主语；第二句的第二个分句没有主语，实际是承上省略了第一个分句的主语“我们”，是个省略的主谓句；第三句的第一个分句没有主语，是个非主谓句；第四句的两个分句都没有主语，都是非主谓句，属于单句中的无主句。但是，它们的两个分句之间，意义都是关联的，彼此联合起来才能表示一个比较完整的意思；两个分句在结构上相对独立，彼此不充当句子成分；有的分句与分句之间有关联词语，把两个分句联系起来；分句与分句之间都有一定的语音停顿。这几点都是复句的特点，因此它们都是复句。

大家可能已经发现，这几个复句的文字并不比一些复杂单句的文字多。如果在复杂单句和复句同时出现的试卷上，我们怎样准确地辨析它们而不致出错呢？

其实，尽管复句的特点很多，但是我们只要抓准了其中根本的一条就好掌握了。这根本的一条，就是看两个分句在结构上是否相对独立，彼此是否不充当句子成分。如果一个分句不充当另一个分句的句子成分，就是说一个完整的句子至少由两个用逗号分开的分句组成，那就是复句；如果一个分句充当另一分句的句子成分，就是说一个完整的句子最多有一个谓语部分，那就是单句。

据此，我们来辨析下面几个句子看看：

（1）\uuline{我}
$\Vert$
\uline{可以}
[大胆地]
\uline{说}，
\uwave{如果欧阳修生活在今天\\
	的话，那他的《秋声赋》一定会是另外一种内容，另外一\\
	种色调}。

（2）\uuline{梁实秋先生}
$\Vert$ [为了《拓荒者》上称他为“资本家的走狗”]，[就]
\uline{做了}（一篇自云“我不生气”的）
\uwave{文章}。

（3）\uuline{不知道谁是他的主子}，
$\Vert$[\uline{正}]
是（他遇见所有阉人都驯良的）
\uline{原因}，[也就]
\uline{是}(属于所有的资本家的)
\uwave{证据}。

（4）\uline{要团结}，[不]
\uwave{要分裂}。

第一句是个单句，因为它的后半部分“如果……一种色调”，尽管独立出来看是个复句，而在整个句子中，这个复句结构却是充当前面动词谓语“说”的宾语成分，回答“说什么的。”

第二句也是个单句，因为从“为了”到“走狗”是个介词结构，充当句子的状语成分，表示动词谓语“做”的原因，不能把它看作是谓语部分，因此就不能把整个句子当作由两个分句构成的复句。

第三句是复句了，而不是单句。因为“不知道谁是它的主子”是第一分句的主语（注意：由于它已作为第一个分句的主语成分，是个动宾词组，因而就没有资格作为一个单句存在了），“是”是谓语，“原因”是宾语；整个句子的后半部分“也就是……证据”是第二个分句，它的主语“不知道谁是它的主子”承前一分句省略了，“是”是第二分句的谓语，“证据”是第二分句的宾语。

第四句也是复句，因为“要团结”和“不要分裂”是两个相对独立的非主谓句，彼此不作句子成分，整个句子有两个谓语部分，全句的意思是由这两个分句联合起来表达的，因此它只能是复句，而不是单句，当然更不能不承认它是一个表达了完整意思的句子。

根据上面的辨析，我们可以这样说：

1. 一个以复句形式出现的复杂词组如果充当了句子成分时（单句作为主谓词组、动宾词组、述补词组等充当句子成分，当然更不在话下），它就失去了句子的资格，不要把它误定为复句。

2. 介词结构既不能充当主语，也不能充当谓语或宾语，只能充当状语、补语或定语，因此不要由于句子中有介词结构而把单句误定为复句。

3. 承前或蒙后省略了主要成分的分句，不能把它误当作别的分句的句子成分，而要把它看作是构成复句的分句。

4. 对于非主谓句中的无主句或独词句，我们不能由于它没有主语或谓语而误把它当作别的分句的句子成分，从而误把复句当成了单句。

至于关联词语，的确也是复句结构上的一个特点。有些复句中分句之间的关系，一定要用关联词语才能表示出来；去掉了必要的关联词语，分句与分句之间的关系就显示不出来。例如：

“不是人们的思想意识决定人们的社会存在，而是人们的社会存在决定人们的思想意识。”

在这个复句中，如果把“不是”、“而是”去掉，不仅看不出分句之间的联系，甚至使人难以理解其意思了。不过，有不少复句并没有关联词语，这就说明关联词语并不是复句结构的决定性因素。

要表示分句之间的不同关系，就得用不同的关联词语。例如：“你肯下苦功，学习成绩好”。前后分句如果分别加上“因为······所以······”、“既然······就······”、“如果······就”、“除非······才······”、“只有······才······”等互相搭配的关联词语，就能构成不同类型的复句，表达不同的意思。

不过，我们应该注意：并不是说句子中有了关联词语，这个句子就一定是复句。因为关联词语除了连接分句之外，还可以连接词、词组、句子，甚至段落，所以在辨析一个句子究竟是单句还是复句时，一定要弄清句子中的关联词语究竟是连接词或词组，还是连接分句的。例如：

（1）文武全才的陈毅元帅，可以而且应该称为“儒帅”。

（2）共产党员应该带头挑起振兴中华这付极为艰巨但又无比光荣的重担。

（3）无论谁都不准无故缺勤。

第一句“而且”连接的是两个词，第二句“但又”连接的是两个词组；第三句的“无论”什么也不连接，用在句子里只起强调的作用。因此，上面的这一类句子都是单句，而不是复句。

另外，复句的分句之间都有逗号（或分号）隔开，表示分句之间的停顿，但是这不等于说凡有逗号隔开的句子都是复句。例如上面第一句中“文武全才的陈毅元帅”之后有逗号，那是为了突出主语而用的句中停顿，决不能认为是一个复句的分句之间的停顿。

总之，有无关联词语或者句中有无逗号，都能帮助我们区分复句和复杂单句，但是都不是最根本的条件。最根本的条件，就是看两个分句在结构上是否彼此不充当句子成分。我们如果抓住这根本的一条，并借助关联词语和逗号，就能得心应手地辨析复句和复杂单句了。

例如：

（1）放下书包，他拿起碗就盛饭。

（这是复句，因为“放下书包”蒙后省略了主语“他”，是个双部句中的省略句；两个分句相对独立，彼此不充当句子成分；分句之间有逗号隔开。）

（2）他从心里感到，没有共产党的领导，就没有劳动人民的彻底解放。

（这是单句，因为“没有共产党的领导，就没有劳动人民的彻底解放”在整个句子里充当了谓语“感到”的宾语成分，全句的主语是“他”。）

（3）没有共产党的领导，就没有劳动人民的解放。

（这是复句，因为用逗号隔开的两个分句都是单部句中的无主句，而且二者的意义相对独立，合起来构成一个完整的意思，特别是彼此不充当句子成分。）

（4）这么大的雨，他恐怕不来了。

（这是复句，因为“这么大的雨”是单部句中的独词句，与后一分句互不充当句子成分，中间又有逗号隔开。）

（5）当革命遇到困难，遭到暂时的失败，处于紧急关头的时候，周总理总是信心百倍，毫不退缩，顽强战斗。

（这是复句，因为“毫不退缩”和“顽强战斗”两个分句都承前省略了主语“周总理”，所以这是由三个并列分句组成的复句，各分句互不充当句子成分，分句之间有逗号隔开。至于“总是”在这个句子中是加强语气的副词，作状语用，并不是判断词。）

（6）争取领导上的帮助和争取群众的批评监督，这是搞好工作的关键。

（这是单句，因为“争取领导上的帮助和争取群众的批评监督”，跟句中的“这”指同一事物，这样的用于句外而跟句中的一个或几个词语指同一事物的成分，叫做外位成分，当然不能作为分句看待。在本句中，这个外位成分与“这”一样，都充当全句的主语成分。）

（7）资产阶级领导的旧民主主义革命无力解救中国，已经被事实证明了。

（这是单句，因为“资产阶级领导的旧民主主义革命无力解救中国”虽然单独拿出来看是个单句，但在这整个句子里却已充当了谓语“证明”的主语成分，所以它已失去了作为单句的资格，只能作为一个主谓词组看。）

（8）真的猛士，敢于直面惨淡的人生，敢于正视淋漓的鲜血。

（这是复句。第二个分句“敢于正视淋漓的鲜血”，承前省略了主语“真的猛士”，所以是个双部分句中的省略句，而且与前一分句互不充当句子成分，又有逗号加以隔开。至于“真的猛士”之后用逗号隔开，只是表示强调主语所作的句中停顿，而不能认为“真的猛士”是个分句。）

（9）我们要把坚持社会主义道路，有专业知识，有组织领导才能，又是年富力强的各种人才，大胆地、有计划地选拔到党的、政府的、经济事业的、科学文化教育事业的领导岗位上来。

（这是单句，主语是“我们”，谓语是“要”“选拔”，“选拔”之前的“把”字句是状语，“选拔”之后的文字是补语。文字虽多，逗号也多，但结构却较简单。）

（10）对于国内外、省内外、县内外、区内外的具体情况，不愿作系统的周密的调查和研究，仅仅根据一知半解，根据“想当然”，就在那里发号施令，这种主观主义的作风，不是还在许多同志中间存在着吗？

（这是单句，因为“对于国内外……就在那里发号施令”与“这种主观主义的作风”是指同一事物，是句子的外位成分，它与“这种主观主义的作风”同作本句的主语成分。）
\newpage